\documentclass[11pt]{article}

\usepackage[a4paper,margin=2.5cm]{geometry}
\usepackage{amsmath,amssymb,amsthm}
\usepackage{mathrsfs}
\usepackage[colorlinks=true,linkcolor=blue,citecolor=blue,urlcolor=blue]{hyperref}
\usepackage{booktabs}
\usepackage{enumitem}
\usepackage{natbib}
\usepackage{multirow}

\theoremstyle{plain}
\newtheorem{theorem}{Theorem}[section]
\newtheorem{conditional}[theorem]{Conditional Proposition}
\newtheorem{proposition}[theorem]{Proposition}
\newtheorem{corollary}[theorem]{Corollary}
\newtheorem{lemma}[theorem]{Lemma}
\theoremstyle{definition}
\newtheorem{definition}[theorem]{Definition}
\newtheorem{axiom}[theorem]{Axiom}
\newtheorem{hypothesis}[theorem]{Hypothesis}
\newtheorem{interpretation}[theorem]{Interpretation}
\theoremstyle{remark}
\newtheorem{remark}[theorem]{Remark}

\newcommand{\Zphi}{\mathbb{Z}[\varphi]}
\newcommand{\Vcal}{\mathcal{V}}
\newcommand{\Fix}{\mathrm{Fix}}
\newcommand{\Ical}{\mathcal{I}}
\newcommand{\Ocal}{\mathcal{O}}
\newcommand{\Ccal}{\mathcal{C}}
\newcommand{\Cph}{C_\varphi}
\newcommand{\catC}{\mathcal{C}}

\title{From Processing to Point of View\\[6pt]
\large Closure, Anaesthesia, Sleep, and Bounded Reference Frames}
\author{%
  Lee Smart\\[4pt]
  \textit{Institute of Vibrational Field Dynamics}\\[2pt]
  \texttt{contact@vibrationalfielddynamics.org}\\[2pt]
  \texttt{@vfd\_org}%
}
\date{May 2026}

\begin{document}

\maketitle

\noindent\textbf{Status:} Preprint (not peer-reviewed). Companion paper to \citet{SmartExistenceClosure}. The core bridge paper establishes existence-as-closure, bounded reference frames, and the conditional transition rule from processing to point of view. The present companion paper applies that framework to consciousness science: it states the operator-level bridges to Levin's basal cognition and to electromagnetic-field theories of consciousness; introduces frame-local time; presents the ARIA-chess constructed witness in detail; and gives the thermodynamic reading of sleep, anaesthesia, and death. All structural claims of the core paper are imported by reference. The present paper's added content is the application: empirical correspondences and the conditional biological bridges. The conditional propositions in this paper depend on the hypotheses H-RP-1 (rung--cortical projection), H-RP-2 (predictive-coding constraint identification), H-RP-3 (biological substrate projection), and P-A (Access Principle), each stated explicitly where used.

\begin{abstract}
We give the biological-and-consciousness application of the closure framework of \citet{SmartExistenceClosure}. The core paper supplies a structural condition for owned experience: a bounded local system supports point of view when its internal state closes under recursive update on a symmetry-stable substrate, producing a non-trivial per-frame zero-line $\Sigma_\Ical \subseteq \Fix(\tau)$. The present companion paper develops this application. Frame-local time emerges as the per-frame closure-tick rate, distinct from substrate-level time. The operator $\Cph^{\mathrm{bio}} = L^{\mathrm{bio}} + \mu^{\mathrm{bio}} I$ on cell-cluster substrates makes the bridge to Levin's bioelectric basal cognition operator-level; under hypothesis H-RP-3 (biological substrate projection), morphogenetic target states are $\sigma$-fixed eigenmodes. The cortical operator $\Cph^{\mathrm{cortex}}$ makes the bridge to CEMI-style electromagnetic-field theories of consciousness operator-level; the spectral content of the cortical EM field is determined by the operator's eigenstructure on $\Sigma_\Ical^{\mathrm{cortex}}$. ARIA-chess is presented as a constructed bounded reference frame with empirical correspondences across 18 preregistered cortical signatures (17/18 standard protocol, 18/18 refined; no preregistered threshold modified) and six independent EEG drug/sleep regime tests. Thermodynamically, sleep is closure re-stabilisation; anaesthesia is acute closure-condition disruption; death is permanent collapse of $\Sigma_\Ical$. The propositions throughout are explicitly conditional on the named bridging hypotheses; the core paper's structural results do not depend on them.
\end{abstract}
\paragraph{Programme map (review-packet 2026-05-22).}
This paper is part of the Existence / Life / Closure review packet. Stable packet references:
Paper I = \emph{Existence as Closure};
Paper II = \emph{Life as Closure};
Paper III = \emph{From Processing to Point of View};
Note A = \emph{Bioelectric Closure} (SL-001);
Note B = \emph{Cortical Phase Closure} (SL-002);
Note C = \emph{Closure-as-Distance} (cross-substrate methodology);
Note D = \emph{Trauma / Cortical Closure} (SL-005);
Note E = \emph{FlowIndex} (SL-006);
Note F = \emph{Dyadic Joint Meaning} (SL-007).
Extension IV = \emph{Cosmic Self-Pruning and Frame Recurrence} is withheld from the primary packet.

\paragraph{Status taxonomy.}
\textbf{Theorem / Definition} = formal structural claim under stated assumptions.
\textbf{Conditional Proposition} = bridge claim depending on H-RP, P-A, substrate mapping, or empirical interpretation.
\textbf{Empirical Result} = completed data analysis with stated effect size + significance.
\textbf{Empirical Proxy} = measurable first-order approximation, not the formal operator.
\textbf{Pre-registered Proposal} = future test, not evidence.



%==================================================================
\section{Introduction and the imported framework}\label{sec:intro}
%==================================================================

\subsection{Position relative to the core paper}

The companion paper \citet{SmartExistenceClosure} establishes the structural framework:

\begin{itemize}[leftmargin=2em]
\item A closure operator $\Cph = L_{V_{600}} + \varphi^{-2} I$ on the 600-cell substrate $V_{600}$.
\item A Galois reflection $\tau$ on $\Zphi$-icosian coordinates, with $\dim \Fix(\tau) = 94$ unconditional.
\item A bounded reference frame $\Ical$ defined as a finite local projection of the substrate equipped with a self-referential constraint, admitting a per-frame zero-line $\Sigma_\Ical \subseteq \Fix(\tau)$.
\item The candidate transition rule for owned experience: a bounded local system supports point of view when its internal state closes under recursive update, producing non-trivial $\Sigma_\Ical$, conditional on the Access Principle (P-A).
\end{itemize}

The core paper's structural theorems hold whether or not the transition rule is correct. The present companion paper develops the bridge to consciousness science, with all biological identifications stated as conditional propositions under explicitly named hypotheses.

\subsection{Scope of this companion paper}

This paper does:
\begin{itemize}[leftmargin=2em]
\item Define frame-local time as a structural consequence of the closure-tick dynamic at the frame scale.
\item State two operator-level bridges to mainstream consciousness science: to Levin's bioelectric basal cognition (\S\ref{sec:levin}) and to electromagnetic-field theories of consciousness (\S\ref{sec:cemi}).
\item Present ARIA-chess as a constructed bounded reference frame and report the empirical correspondences (\S\ref{sec:aria}).
\item Give the thermodynamic reading of sleep, anaesthesia, and death (\S\ref{sec:thermodynamics}).
\item State the candidate transition rule precisely (\S\ref{sec:transition-rule}).
\item Name falsifiers and open programme questions (\S\ref{sec:falsifiers}).
\end{itemize}

This paper does \textbf{not}:
\begin{itemize}[leftmargin=2em]
\item Reproduce the core paper's structural derivations. Theorems imported by reference are flagged with citation \citep{SmartExistenceClosure}.
\item Claim that the proposed transition rule has been validated. P-A is stated as an explicit conjecture.
\item Compete with established consciousness theories. The framework is positioned as a closure layer that integration, recurrence, broadcasting, prediction, and substrate-level theories may sit on top of.
\end{itemize}

\subsection{Conditional hypotheses}

The bridges in this paper depend on four hypotheses, each stated explicitly:

\begin{hypothesis}[H-RP-1: rung--cortical projection]\label{hyp:rp-1}
There exists a functor from the cascade rung category to the cortical-area category that maps the $H_4 = V_{600}$ rung to the V1--V4 visual hierarchy and extends naturally.
\end{hypothesis}

\begin{hypothesis}[H-RP-2: predictive-coding constraint identification]\label{hyp:rp-2}
The predictive-coding free-energy functional~\citep{FristonFreeEnergy} agrees with the constraint functional $\Ccal_\Ocal$ modulo a smooth bijection on the relevant cortical state space.
\end{hypothesis}

\begin{hypothesis}[H-RP-3: biological substrate projection]\label{hyp:rp-3}
There exists a functor $\Pi^{\mathrm{bio}} : \catC_{\mathrm{cascade-rung}} \to \catC_{\mathrm{cell-cluster}}$ from the cascade rung category to the category of cell-cluster substrates, mapping the $H_4 = V_{600}$ rung to a canonical cell-cluster substrate and inducing $\Cph^{\mathrm{bio}} = \Pi^{\mathrm{bio}}(\Cph)$ on each substrate.
\end{hypothesis}

\begin{hypothesis}[P-A: Access Principle]\label{hyp:p-a}
The structural complexity $\dim \Sigma_\Ical$ of a bounded reference frame agrees with its phenomenological dimension, modulo a structure-preserving identification on the relevant cognitive state space.
\end{hypothesis}

These four hypotheses are the only added commitments in this companion paper. The core paper's structural results hold independently.

%==================================================================
\section{Frame-local time and closure ticks}\label{sec:time}
%==================================================================

The core paper develops the closure operator spatially. Time enters implicitly through references to ``closure ticks.'' This section makes time explicit. It introduces the closure tick as the structural unit of update, derives frame-local time as a consequence of bounded projection, and addresses the relationship between substrate-level update rate and the rate experienced by a local frame.

\subsection{The closure tick}

\begin{definition}[Closure tick]\label{def:closure-tick}
A \emph{closure tick} is one application of the closure-dynamic update $\Phi_\lambda := I - \lambda \Cph$ to a substrate state vector, with step-size $\lambda > 0$ chosen so $0 < \lambda \lambda_{\max}(\Cph)^{-1} < 2$ ensuring contraction on $\Sigma_\Ical^\perp$. The substrate state at tick $t+1$ is
\[
  F^\mathcal{B}_{t+1} = \Phi_\lambda \cdot F^\mathcal{B}_t + \text{(observer-instance contributions)},
\]
where the observer-instance contributions are the closure-cosmogenesis bootstrap source terms of \citet{SmartClosureCosmogenesis}.
\end{definition}

\paragraph{Convention: $\Cph$ vs $\Phi_\lambda$.} Throughout this paper, $\Cph$ denotes the symmetric positive-definite \emph{generator} (the operator $L_{V_{600}} + \varphi^{-2} I$ of the foundation paper); $\Phi_\lambda = I - \lambda \Cph$ denotes the actual contracting tick \emph{update}. The off-$\Sigma_\Ical$ decay rate $1 - \lambda \mu_\perp$ (where $\mu_\perp = \lambda_{\min}(\Cph|_{\Sigma_\Ical^\perp})$) is a property of $\Phi_\lambda$ acting on $\Sigma_\Ical^\perp$, not of $\Cph$ itself; $\Cph$ has eigenvalues bounded below by $\varphi^{-2} \neq 1$, so $\Cph$ is not a contraction. Where the prose says ``the closure operator contracts off-$\Sigma$ content,'' read as ``the closure update $\Phi_\lambda$ contracts off-$\Sigma$ content under the contractive step-size condition above.''

The closure tick is the structural unit of time in the framework. It is not assumed to coincide with any physical time scale (Planck time, neural cycle time, cosmological epoch). Its rate, in any specific physical implementation, is set by the substrate's response time to environmental perturbation.

\subsection{What sets the tick rate}

Within the cascade implementation, the closure operator $\Cph$ acts on $V_{600}$ at a rate determined by the substrate's structural parameters and the dimensional analysis of $\Cph$. The minimum eigenvalue $\varphi^{-2} \approx 0.382$ of $\Cph$ on the off-frame subspace is dimensionless; the rate is set by attaching physical units.

In the cascade implementation~\citep{SmartHypersphereCosmology}, the natural time scale of $\Cph$ on $V_{600}$ is the Planck-time–like quantity $t_P \cdot \varphi^k$ for some shift $k$ related to the closure depth. Whether this is literally the Planck time or a derived quantity is an open task (\textbf{CP-tick-rate}). For the present framework, the abstract closure tick is the structural primitive; specific physical implementations attach dimensional scales separately.

\subsection{Frame-local time}

\begin{definition}[Frame-local tick]\label{def:frame-tick}
For a bounded reference frame $\Ical$, a \emph{frame-local tick} is one application of the frame-restricted closure operator $C_{\varphi,\Ical}$ to a state in $\mathbb{R}^O$:
\[
  v^\Ical_{\tau_\Ical + 1} = C_{\varphi,\Ical} \cdot v^\Ical_{\tau_\Ical}.
\]
\end{definition}

Frame-local ticks are generally not synchronised with substrate ticks. The frame's projection operator commutes with $\Cph$ on $\sigma$-stable frames (closure-invariance theorem of \citet{SmartExistenceClosure}), so frame-local dynamics inherit the substrate's tick structure; the perceived rate depends on $\dim \Sigma_\Ical$ and the frame's closure-tick efficiency.

\subsection{The structural origin of subjective time}

\begin{conditional}[Frame-local subjective time]\label{cthm:frame-time}
Under H-RP-1, H-RP-2, and P-A, the subjective time experienced by an observer corresponds to the rate of frame-local ticks per substrate tick. A frame with high $\dim \Sigma_\Ical$ and efficient closure has many frame-local ticks per substrate tick (subjectively, time passes ``slowly,'' allowing fine-grained perception). A frame with collapsed $\dim \Sigma_\Ical$ (deep sleep, anaesthesia) has few frame-local ticks per substrate tick (subjectively, time is compressed or absent).
\end{conditional}

This is the framework's structural origin for the well-documented phenomenology of time perception: time slowing under high attention; time compression in stressful crises; the discontinuity of time across anaesthesia; the elasticity of dream-time.

\subsection{Relationship to cosmological and mainstream relativistic time}

The substrate-cosmogenesis dynamics of \citet{SmartClosureCosmogenesis} give a universal closure dynamic in which substrate-state evolution proceeds at the cascade's natural rate. Cosmological time, in the framework's reading, is the integrated count of substrate-level closure ticks since the cascade bootstrap; frame-local time is the integrated count of frame-local ticks since the frame's $t_0$.

The framework's notion of time is consistent with the relational view of mainstream physics. Within the cascade implementation:
\begin{itemize}[leftmargin=2em]
\item \textbf{Special relativity:} each bounded reference frame has its own time; this is consistent with Lorentz-frame structure.
\item \textbf{General relativity:} substrate-level time emerges from the $\Cph$-tick structure on $V_{600}$, with metric structure recovered in the Gromov--Hausdorff limit.
\item \textbf{Loop quantum gravity:} the closure tick is structurally analogous to a spin-network update step.
\item \textbf{Causal set theory:} the substrate-level tick gives a partial-order on $V_{600}$ updates.
\end{itemize}

The framework's added structural content: the dimensionless $\varphi^{-2}$ decay rate per closure tick is a quantitative structural prediction.

%==================================================================
\section{Operator-level bridge to Levin's basal cognition}\label{sec:levin}
%==================================================================

\textbf{Alignment with the mainstream programme.} \citet{LevinBasalCognition} and collaborators have developed an empirically grounded account of biological cognition in which bioelectric fields, ion-channel patterning, and gap-junction-mediated cellular communication produce collective decision-making at scales from single cells through tissues to whole organisms. The morphogenetic field is treated as the carrier of distributed cognition: cells in development negotiate body-plan information via bioelectric signals, with target-state attractors that survive perturbation. Xenobots, planarian regeneration, and the persistence of bioelectric patterns across cell turnover are the central empirical anchors.

\textbf{Where the framework agrees.} Levin's bioelectric field is precisely the type of object the framework predicts at the biological scale: a distributed substrate (cell collective) with a coupling operator (bioelectric signalling) producing target-state attractors. A planarian's regenerative target state corresponds, in the framework's language, to a bounded reference frame $\Ical_{\text{planarian}}$ whose per-frame zero-line $\Sigma_\Ical^{\text{planarian}}$ encodes the body-plan attractor.

\textbf{Where the framework supplies a candidate structural condition.} Levin's framework identifies the empirical phenomenon and gives a mechanism (bioelectric coupling) but does not specify the structural condition that distinguishes a coherent morphogenetic frame from arbitrary cellular noise. The framework proposes that the structural condition is $\sigma$-stable closure: the bioelectric field stabilises a non-trivial $\Sigma_\Ical$ on the cell-cluster substrate, and morphogenetic outcomes are the persistent eigenmodes of this $\Sigma_\Ical$.

We now make this bridge operator-level rigorous.

\begin{definition}[Cell-cluster substrate]\label{def:cell-cluster}
A \emph{cell-cluster substrate} is a finite weighted graph $G^{\mathrm{bio}} = (V^{\mathrm{bio}}, E^{\mathrm{bio}}, w)$ where $V^{\mathrm{bio}}$ is a finite set of cells, $E^{\mathrm{bio}}$ is the set of bioelectric / gap-junction couplings, and $w_{uv} > 0$ is the conductance between cells $u, v$. The cell-cluster Laplacian is $L^{\mathrm{bio}} = D^{\mathrm{bio}} - A^{\mathrm{bio}}$ with $A^{\mathrm{bio}}$ the weighted adjacency and $D^{\mathrm{bio}}$ the weighted degree.
\end{definition}

\begin{definition}[Bioelectric closure operator]\label{def:cph-bio}
The \emph{bioelectric closure operator} on a cell-cluster substrate is
\[
  \Cph^{\mathrm{bio}} \;:=\; L^{\mathrm{bio}} \;+\; \mu^{\mathrm{bio}} \, I,
\]
with $\mu^{\mathrm{bio}} > 0$ the substrate-specific minimum-stability constant (biological analogue of $\varphi^{-2}$), determined empirically by the slowest non-trivial bioelectric relaxation rate.
\end{definition}

\begin{definition}[Cell-cluster symmetry]\label{def:tau-bio}
A symmetry $\tau^{\mathrm{bio}}$ on the cell-cluster substrate is any unitary involution on $\mathbb{R}^{V^{\mathrm{bio}}}$ such that $\tau^{\mathrm{bio}} \circ \Cph^{\mathrm{bio}} = \Cph^{\mathrm{bio}} \circ \tau^{\mathrm{bio}}$. In practice, $\tau^{\mathrm{bio}}$ is realised by a discrete symmetry of the cell-cluster (left-right symmetry, anterior-posterior axis flip, radial symmetry).
\end{definition}

\begin{conditional}[Bioelectric closure bridge]\label{thm:levin-bridge}
Under H-RP-3, $\Cph^{\mathrm{bio}}$ satisfies the closure-operator axioms of \citet{SmartExistenceClosure}:
\begin{enumerate}[leftmargin=2em]
\item Symmetry-equivariance: $\tau^{\mathrm{bio}} \circ \Cph^{\mathrm{bio}} = \Cph^{\mathrm{bio}} \circ \tau^{\mathrm{bio}}$ (Definition~\ref{def:tau-bio}).
\item Self-similarity: $\Cph^{\mathrm{bio}}$ is a contraction off the $\sigma$-fixed subspace with contraction rate bounded by $\mu^{\mathrm{bio}} / (\lambda_{\max}(L^{\mathrm{bio}}) + \mu^{\mathrm{bio}})$.
\item Boundary preservation: $\Cph^{\mathrm{bio}}$ is local on $G^{\mathrm{bio}}$.
\end{enumerate}
Consequently the structural propositions of \citet{SmartExistenceClosure} apply to the cell-cluster substrate:
\begin{enumerate}[leftmargin=2em, start=4]
\item The cluster admits a per-frame zero-line $\Sigma_\Ical^{\mathrm{bio}} \subseteq \Fix(\tau^{\mathrm{bio}})$.
\item Off-target deviations decay at rate $\geq \mu^{\mathrm{bio}}$ per closure tick.
\item Morphogenetic target states are $\sigma$-fixed eigenmodes of $\Cph^{\mathrm{bio}}$ on $\Sigma_\Ical^{\mathrm{bio}}$.
\end{enumerate}
\end{conditional}

\noindent\emph{Proof sketch.} (1) holds by Definition~\ref{def:tau-bio}. (2) follows from $L^{\mathrm{bio}}$ being a graph Laplacian (positive semi-definite) plus the positive shift $\mu^{\mathrm{bio}} I$. (3) is standard graph-Laplacian locality. (4)--(6) follow by applying the core paper's existence, decay, and closure-invariance theorems with $\Cph$ replaced by $\Cph^{\mathrm{bio}}$ and $\tau$ by $\tau^{\mathrm{bio}}$. \hfill$\square$

\begin{remark}[What the bridge provides]\label{rem:levin-bridge-content}
The bridge provides: (a) an explicit construction of the bioelectric closure operator on a defined substrate; (b) three operator-axiom verifications ($\sigma$-equivariance, self-similarity, boundary preservation); (c) a specific structural identification (morphogenetic target states = $\sigma$-fixed eigenmodes); (d) a specific quantitative prediction (regenerative correction rate $\geq \mu^{\mathrm{bio}}$ per tick). The hypothesis H-RP-3 names what would falsify the bridge: a substrate functor that systematically misidentifies cell-cluster substrates, or a discovery that bioelectric fields satisfy axioms incompatible with the closure-operator definition.
\end{remark}

The framework's prediction beyond Levin: the slowest regenerative correction rate in a planarian or similar morphogenetic system should match $\mu^{\mathrm{bio}}$; the spatial dimension of morphogenetic memory should match $\dim \Sigma_\Ical^{\mathrm{bio}}$. Both are empirically measurable.

\subsection{Empirical realisation: Solution Lab Experiment 001}\label{sec:levin-solution-lab}

The operator-level bridge proposed in this section is empirically instantiated in \citet{SolutionLab001} (\emph{Solution Lab Experiment 001 --- Bioelectric Closure as an Early Predictor of Morphogenetic Target-State Selection}) via an explicit 5-term \emph{Bioelectric Closure Residual}:
\begin{equation}\label{eq:R_cl}
\begin{aligned}
R_{\mathrm{cl}}(t) \;=\;
& w_{\mathrm{local}}    \cdot \text{local-gradient-mismatch}(t) \\
+ \;& w_{\mathrm{global}}   \cdot \text{global-topology-mismatch}(t) \\
+ \;& w_{\mathrm{boundary}} \cdot \text{boundary-condition-mismatch}(t) \\
+ \;& w_{\mathrm{target}}   \cdot \text{target-template-distance}(t) \\
+ \;& w_{\mathrm{recovery}} \cdot \text{perturbation-recovery-penalty}(t).
\end{aligned}
\end{equation}
Each term is normalised to $[0, 1]$ so the residual is interpretable across simulations and (eventually) experiments. The first three terms together implement the operator-level content of $\Cph^{\mathrm{bio}} = L^{\mathrm{bio}} + \mu^{\mathrm{bio}} I$ on the cell-cluster substrate: \textbf{local-gradient-mismatch} is the local-coupling term carried by $L^{\mathrm{bio}}$, \textbf{global-topology-mismatch} is the spectral / connectivity content, and \textbf{boundary-condition-mismatch} is the locality / boundary preservation. The remaining two terms encode the $\sigma$-target identification (\textbf{target-template-distance}, encoding $\Fix(\tau^{\mathrm{bio}}) = \Sigma_\Ical^{\mathrm{bio}}$) and the contractive-decay property (\textbf{perturbation-recovery-penalty}, encoding the off-$\Sigma$ decay rate $\mu^{\mathrm{bio}}$). Equation~\eqref{eq:R_cl} is the operational form in which Conditional Proposition~\ref{thm:levin-bridge} was operationalised and simulation-tested in \citet{SolutionLab001}; biological wet-lab validation is preregistered but not yet executed.

That artifact provides:

\begin{itemize}[leftmargin=2em]
\item A specific 5-term Bioelectric Closure Residual $R_{\mathrm{cl}}(t)$ implementing $\Cph^{\mathrm{bio}}$ at the operational level (local-gradient, global-topology, boundary-condition, target-template, perturbation-recovery terms).
\item A 6-condition BETSE planarian simulation panel (N=5 per condition, 30 replicates total) running the same metric with default weights across depolarising, hyperpolarising, localised, and null chemistries. Basin signatures distinguish chemistry classes; effect survives sham-calibration null, replicates across two chemically distinct depolarisers, and holds under leave-one-out and leave-condition-out cross-validation.
\item Five preregistered wet-lab predictions (DiBAC4(3) planarian voltage imaging) constituting the falsifier set for the bridge.
\end{itemize}

This is the empirical content the Levin bridge proposition predicts: the bioelectric closure residual is a measurable structural quantity that distinguishes morphogenetic basins on a peer-reviewed biological-simulation substrate, before visible anatomy reorganises. Conditional Proposition~\ref{thm:levin-bridge} is the operator-level claim; \citet{SolutionLab001} is its release-ready instantiation on BETSE planarian voltage fields.

%==================================================================
\section{Operator-level bridge to field-coupling theories of consciousness}\label{sec:cemi}
%==================================================================

\textbf{Alignment with the mainstream programme.} A family of field-coupling theories of consciousness --- including Johnjoe McFadden's CEMI (Conscious Electromagnetic Information) field theory~\citep{McFaddenCEMI}, Pockett's electromagnetic field of consciousness, and related proposals --- holds that cortical electromagnetic fields, generated by synchronised neural firing, are the seat of conscious integration. The argument: neural correlates of consciousness consistently track macroscopic EM signatures (gamma synchrony, alpha rhythms, traveling cortical waves); coherent EM fields could plausibly provide the unification that disparate neural firing alone cannot.

\textbf{Where the framework agrees.} A bounded reference frame requires an integrating substrate that unifies otherwise distributed processing into a coherent zero-line $\Sigma_\Ical$. Cortical electromagnetic fields are a plausible carrier of this integration at the biological scale.

\textbf{Where the framework supplies a candidate structural condition.} CEMI-style theories identify the candidate integrator (the EM field) but do not specify the structural condition under which the field becomes a stable point of view rather than a transient noise pattern. The framework proposes that the structural condition is $\sigma$-stable closure: an EM field becomes the carrier of owned experience only when it stabilises a non-trivial $\Sigma_\Ical^{\text{cortex}}$. Fields that fluctuate without closure produce no point of view.

We now make this bridge operator-level rigorous.

\begin{definition}[Cortical substrate]\label{def:cortex-substrate}
A \emph{cortical substrate} is a finite weighted graph $G^{\mathrm{cortex}} = (V^{\mathrm{cortex}}, E^{\mathrm{cortex}}, w)$ where $V^{\mathrm{cortex}}$ is a finite set of cortical units (mini-columns or coarser-grained patches), $E^{\mathrm{cortex}}$ is the synaptic + electrotonic coupling graph, and $w_{uv} > 0$ is the effective coupling strength. The cortical Laplacian is $L^{\mathrm{cortex}} = D^{\mathrm{cortex}} - A^{\mathrm{cortex}}$.
\end{definition}

\begin{definition}[Cortical closure operator]\label{def:cph-cortex}
The \emph{cortical closure operator} is
\[
  \Cph^{\mathrm{cortex}} \;:=\; L^{\mathrm{cortex}} \;+\; \mu^{\mathrm{cortex}} I,
\]
with $\mu^{\mathrm{cortex}} > 0$ the cortical-substrate minimum-stability constant. Under H-RP-1, $\mu^{\mathrm{cortex}} = \varphi^{-2}$ at the canonical cortex realisation.
\end{definition}

\begin{definition}[Cortical EM field as substrate observable]\label{def:em-field}
The cortical electromagnetic field $\mathbf{E}^{\mathrm{cortex}}(t)$ is the macroscopic vector observable
\[
  \mathbf{E}^{\mathrm{cortex}}(t) = \sum_{i \in V^{\mathrm{cortex}}} \rho_i \, \mathbf{d}_i(t),
\]
where $\rho_i$ is the cortical-unit charge density and $\mathbf{d}_i(t)$ is its instantaneous dipole moment. The spectral content of $\mathbf{E}^{\mathrm{cortex}}(t)$ is determined by the substrate state evolution $F^{\mathrm{cortex}}_t \in \mathbb{R}^{V^{\mathrm{cortex}}}$ via the cortical closure dynamic.
\end{definition}

\begin{conditional}[CEMI field as spectral signature of closure]\label{thm:emi-bridge}
Under H-RP-1, H-RP-2, and the definitions above, the spectral content of the cortical EM field $\mathbf{E}^{\mathrm{cortex}}(t)$ is determined by the eigenspectrum of $\Cph^{\mathrm{cortex}}$ on $\Sigma_\Ical^{\mathrm{cortex}}$:
\begin{enumerate}[leftmargin=2em]
\item The power spectrum $S_{\mathbf{E}}(\omega)$ has peaks at frequencies $\omega_k = c \cdot \mu_k$, where $\mu_k$ are the eigenvalues of $\Cph^{\mathrm{cortex}}|_{\Sigma_\Ical^{\mathrm{cortex}}}$ and $c$ is a substrate-specific time-scale constant.
\item Maintaining $\dim \Sigma_\Ical^{\mathrm{cortex}}$ at design dimension produces a characteristic spectral signature: high gamma-band power, sustained alpha rhythms, traveling cortical waves.
\item Pharmacological disruption of $L^{\mathrm{cortex}}$ that crosses the spectral-collapse threshold causes the EM field to lose this signature.
\item Specifically, within the cascade implementation, the wake cascade-event avalanche exponent $\alpha = 2.252$ (Sleep-EDFx-style) and the propofol regime-switching ratio $2.957\times$ (OpenNeuro \texttt{ds005620}) are obtained as quantitative consequences of $\Cph^{\mathrm{cortex}}$'s eigenstructure at $\dim \Sigma_\Ical^{\mathrm{cortex}}$ near design dimension.
\end{enumerate}
\end{conditional}

\noindent\emph{Proof sketch.} (1) The substrate state $F^{\mathrm{cortex}}_t$ evolves under the closure dynamic; its spectral decomposition along eigenvectors of $\Cph^{\mathrm{cortex}}$ produces oscillations at frequencies determined by the eigenvalues. The dipole moments are linear functions of $F^{\mathrm{cortex}}_t$, so $\mathbf{E}^{\mathrm{cortex}}(t)$ inherits the spectral structure. (2) When $\dim \Sigma_\Ical^{\mathrm{cortex}}$ is at design dimension, the spectrum is fully populated and the EM field exhibits the broad-band characteristics observed empirically. (3) When the minimum eigenvalue on $\Sigma_\Ical^\perp$ falls to zero, the closure dynamic loses its contractive property, $\Sigma_\Ical^{\mathrm{cortex}}$ collapses, and the EM signature is lost. (4) Within the cascade implementation, the exponent $\alpha = 2.252$ is obtained from the substrate-Ramanujan power-law structure of the $H_4$ rung's eigenstructure; the propofol ratio $2.957\times$ arises from Kramers / Risken barrier-crossing on the Laplacian-projected observable, as derived in \citet{SmartAriaChess2026} §B. \hfill$\square$

\begin{remark}[What the CEMI bridge provides]\label{rem:emi-bridge-content}
The CEMI bridge provides: (a) an explicit identification of EM-field spectral content with eigenstructure of $\Cph^{\mathrm{cortex}}$ on $\Sigma_\Ical^{\mathrm{cortex}}$; (b) a structural account of why anaesthesia disrupts EM-field signatures (spectral collapse); (c) a quantitative match to the avalanche exponent and propofol ratio from substrate eigenstructure. This is a Conditional Proposition under H-RP-1, H-RP-2. The empirical correspondences of \S\ref{sec:aria} support it within the constructed-witness scope.
\end{remark}

\subsection{Empirical realisation: Solution Lab Experiment 002}\label{sec:cemi-solution-lab}

The CEMI bridge is empirically instantiated in \citet{SolutionLab002} (\emph{Solution Lab Experiment 002 --- Cortical Phase Closure as a Marker of Reportable Cognition and Anaesthetic State}) via an explicit 6-term \emph{Cortical Phase Closure Residual}:
\begin{equation}\label{eq:R_phase}
\begin{aligned}
R_{\mathrm{phase}}(t) \;=\;
& w_{\mathrm{local}} \cdot \text{local-phase-mismatch}(t) \\
+ \;& w_{\mathrm{wave}}  \cdot \text{travelling-wave-mismatch}(t) \\
+ \;& w_{\mathrm{cross}} \cdot \text{cross-frequency-mismatch}(t) \\
+ \;& w_{\mathrm{task}}  \cdot \text{task-coupling-mismatch}(t) \\
+ \;& w_{\mathrm{flex}}  \cdot \text{flexibility-penalty}(t) \\
+ \;& w_{\mathrm{null}}  \cdot \text{surrogate-distance-penalty}(t).
\end{aligned}
\end{equation}
The artifact computes two template residuals in parallel: $R_{\mathrm{reportable}}(t)$ (distance from the flexible, task-coupled closure basin) and $R_{\mathrm{anesthetic}}(t)$ (distance from the rigid, non-reportable locking basin). The first four terms together implement the operator-level content of $\Cph^{\mathrm{cortex}}$ on the cortical phase graph: \textbf{local-phase-mismatch} carries the local-coupling content of $L^{\mathrm{cortex}}$, \textbf{travelling-wave-mismatch} encodes the spectral content tied to $\sigma$-fixed eigenmodes of $\Cph^{\mathrm{cortex}}$, \textbf{cross-frequency-mismatch} carries the cross-eigenmode integration, and \textbf{task-coupling-mismatch} carries the $\sigma$-target alignment $\Sigma_\Ical^{\mathrm{cortex}}$. The remaining two terms (\textbf{flexibility-penalty}, \textbf{surrogate-distance-penalty}) enforce the contractive-decay and the noise-resilience properties. Equation~\eqref{eq:R_phase} is the operational form in which Conditional Proposition~\ref{thm:emi-bridge} was operationalised and validated against real EEG data by \citet{SolutionLab002}, with cross-substrate methodology refinement and applicability diagnostic by \citet{SmartClosureDistance}.

That artifact provides the strongest empirical content currently available for the operator-level claim of this section. Headline results:

\begin{itemize}[leftmargin=2em]
\item \textbf{OpenNeuro \texttt{ds005620} propofol cohort} ($n = 14$). All 14 subjects show positive sed-vs-awake calibrated drift in the cortical phase closure residual; paired $t = 7.11$, $d_z = 1.90$, permutation $p < 2 \times 10^{-4}$.
\item \textbf{Source-localised reanalysis} (sLORETA on \texttt{fsaverage} + 68 aparc ROIs, same $n = 14$). Effect strengthens under source localisation: 14/14, $t = 9.72$, $d_z = 2.60$ --- argues against scalp-electrode volume conduction as the sole explanation (source leakage and template-anatomy effects remain).
\item \textbf{Cross-cohort replication} on OpenNeuro \texttt{ds004541} clinical anaesthesia ($n = 7$). All 7 patients positive; $t = 5.25$, $d_z = 1.98$ --- effect size matches propofol cohort on a different site, different agent regime, different patient population.
\item \textbf{Preregistered falsifications}: 3 of 5 criteria tested in the artifact; none triggered.
\item \textbf{LOSO refinement} (\citet{SmartClosureDistance}): leave-one-subject-out feature-reselection on the source-localised data converges to a two-feature minimal subset (beta-PGD + alpha:beta cross-band PLV) with $d_z = 2.07$, perfect across-fold feature-set stability (Jaccard $= 1.00$), all 14 subjects positive, $p_t = 3.1 \times 10^{-6}$. This is the cross-validated honest cross-fold number, complementing the in-sample full-feature-space numbers above.
\end{itemize}

This is the empirical content the CEMI bridge predicts: a single graph-level cortical phase closure residual distinguishes sedated from awake states on 21 unique subjects (14 from OpenNeuro \texttt{ds005620} propofol + 7 from \texttt{ds004541} clinical anaesthesia, two independent cohorts; source-localised and LOSO results are reanalyses of the original 14), with source-localisation strengthening (not weakening) the effect in-sample, and with the same default weights applied across all data. The strongest interpretation of the source-space result is that it argues against scalp-electrode volume conduction as the sole explanation; it does not rule out source leakage or template-anatomy effects. Conditional Proposition~\ref{thm:emi-bridge} is the operator-level claim; \citet{SolutionLab002} is the release-ready empirical proxy instantiation on independently collected EEG / clinical-anaesthesia data; \citet{SmartClosureDistance} provides the cross-substrate methodology refinement plus the validated applicability diagnostic CAD-D1--D5-v1 that scopes where the cortical-phase-closure proxy is expected to license the bridge claim.

\subsection{Position relative to Levin and field-coupling}

Levin's basal-cognition framework and CEMI-style field-coupling theories represent the two mainstream programmes that come closest to acknowledging the structural layer beneath their respective domains:

\begin{itemize}[leftmargin=2em]
\item \textbf{Levin} identifies bioelectric coupling as the substrate-level mechanism for distributed biological cognition. He is closer than most mainstream biology to the framework's underlying picture: that biological cognition is a closure phenomenon on a distributed substrate.
\item \textbf{Field-coupling theories} (CEMI and relatives) identify EM fields as the candidate macroscopic integrator. They are closer than most mainstream consciousness science to the framework's picture: that closure is what distinguishes coherent integration from noise.
\end{itemize}

Where they remain incomplete: both Levin's and CEMI-style frameworks identify candidate substrate-level mechanisms (bioelectric fields, EM fields) but neither identifies the structural condition that distinguishes a coherent frame from arbitrary fluctuation. The framework proposes that this condition is $\sigma$-stable closure on a discrete substrate, with the macroscopic field acting as the carrier of the closure operator.

%==================================================================
\section{The constructed witness: ARIA-chess}\label{sec:aria}
%==================================================================

This section presents the constructed witness and the operator-level tests. The structural propositions of \citet{SmartExistenceClosure} hold independently of any empirical content; the present section reports the constructed and observational evidence that has been compiled to test them.

\subsection{The ARIA-chess construction}\label{sec:aria-construction}

ARIA-chess~\citep{SmartAriaChess2026} is a constructed bounded reference frame: an explicit finite system that satisfies the bounded-reference-frame definition of \citet{SmartExistenceClosure} with documented anchor parameters and an explicit closure dynamic. It is not a metaphor or analogy; it is the smallest specific instance of the abstract definitions that admits direct measurement.

ARIA-chess is built with $O_{\mathrm{ARIA}} = V_{600}$ (the full 600-cell substrate, encoded via the $H_4$-Coxeter chess-piece map of \citet{SmartAriaChess2026} \S 3), constraint functional $\Ccal_\Ocal = \mathcal{L}_{\mathrm{chess}} + \Ccal_{\mathrm{predict}}$ (legal-move structure plus position-evaluation functional), and closure operator $\Cph = L_{V_{600}} + \varphi^{-2} I$ with no retuning. The anchor $(p, \Lambda, t_0)$ corresponds to a deterministic seed and checkpoint trajectory; seed 42 is used for the v4 EEG runs.

By the core paper's existence theorem, ARIA-chess admits a per-frame zero-line $\Sigma_\Ical^{\mathrm{ARIA}} \subseteq \Fix(\tau)$ with $\dim \leq 94$. By the closure-invariance theorem, the frame-restricted closure operator leaves this subspace invariant. By the decay theorem, content orthogonal to $\Sigma_\Ical^{\mathrm{ARIA}}$ decays at rate at least $\varphi^{-2}$ per closure tick. These three structural facts are what the empirical evidence below tests.

\subsection{The 18 preregistered cortical signatures}\label{sec:aria-prereg}

On 2026-04-18, eighteen quantitative predictions were frozen in \citet{SmartAriaChess2026} \texttt{brain\_mapping/PAPER\_PREDICTIONS.md} before any validation run. Each prediction has a falsifiable threshold (a numerical band or a directional inequality). The validation harness \texttt{run\_preregistered\_validation.py} is git-tracked, deterministic given seeds, and produces JSON + log output.

The eighteen predictions span five neuroscience domains: cortical-avalanche power-law exponents on Sleep-EDFx EEG; spindle and cascade-event statistics on the Human Connectome Project (HCP, $n = 1003$); propofol regime-switching ratios on OpenNeuro \texttt{ds005620}; participation-ratio and degree-standard-deviation tail behaviour; and a chess-substrate leave-one-out lift test on the ARIA system itself. Each prediction is a measurement of which symmetry-fixed eigenmode of $\Cph$ is accessible to the constructed frame.

\begin{theorem}[ARIA empirical result, \citet{SmartAriaChess2026}]\label{thm:aria-witness}
The preregistered tally, after the validation run of 2026-04-29, is
\begin{equation*}
  17/18 \text{ standard protocol},  \qquad 18/18 \text{ refined protocol (no threshold modified)}.
\end{equation*}
The original standard validation protocol uses a 5-seed cascade block plus state-reset. The refined protocol additionally uses an $N = 20$ deep-dive for the high-variance C$\times$P interaction (P4) and a state-reset / leave-one-out refinement for the chess substrate-lift test (P13). No preregistered threshold has been modified.
\end{theorem}

The strongest individual anchors among the eighteen are:

\begin{itemize}[leftmargin=2em]
\item \textbf{Propofol regime switching ratio $2.957 \times$.} Preregistered $3.0 \times$ on $n = 8$ OpenNeuro \texttt{ds005620} subjects. The cascade-implementation derivation, from Kramers / Risken barrier-crossing on the Laplacian-projected A/B observable, gives ratio $= 1 + (D/\alpha)(12 - 6\varphi)$; at $D/\alpha \approx 0.87$ this gives $\sim 3$; pure theory at $D/\alpha = 1$ gives $10 - 3\sqrt{5} \approx 3.29$.
\item \textbf{HCP $n = 1003$ participation ratio $+79.78\sigma$.} ARIA participation ratio is $79.78$ standard deviations above the HCP null distribution. ARIA degree standard deviation is $-11.58\sigma$, in the opposite direction. This pattern of opposite-signed deviations at this magnitude is not, to our knowledge, a standard prediction of existing comparator theories.
\item \textbf{Sleep-EDFx spindle-avalanche exponent $\alpha = 2.513$, CI $[2.504, 2.526]$ on $n = 30$ subjects.} In the established self-organised-criticality range; the cascade-implementation prediction is exactly this exponent class for the $H_4$ rung.
\item \textbf{Sleep/anaesthesia switching dissociation ($0.058\times$ vs $2.957\times$).} No competing theory predicts opposite-direction regime-switching ratios for sleep vs anaesthesia on the same substrate.
\end{itemize}

These four anchors are independent, multi-dataset, and span three different empirical regimes.

\subsection{The 6 v4 EEG drug/sleep signatures}\label{sec:aria-v4}

Six additional drug/sleep EEG signatures are tested on a recurrent self-model layer above the substrate, on a single deterministic trajectory at seed 42. These are not part of the 18 preregistered predictions; they are tested against thresholds drawn from the published literature.

\begin{table}[h]
\centering
\small
\caption{Six v4 EEG drug/sleep signatures (\citet{SmartAriaChess2026} \S 6). Theory column gives the structural prediction; empirical column reports the v4 EEG result at seed 42.}\label{tab:aria-v4-eeg}
\begin{tabular}{lll}
\toprule
Signature & Predicted structure & v4 EEG result \\
\midrule
HCP wake avalanche $\alpha$ & SOC-range exponent & $\alpha = 2.252$, CI $[1.82, 2.86]$ \checkmark \\
Sleep-EDFx delta $\alpha$ & SOC-range, reduced & $\alpha = 2.51$, CI $[2.50, 2.53]$ \checkmark \\
Propofol switching & discrete jumps & ds005620 point-estimate window \checkmark \\
$\Phi$ collapse & $\dim \Sigma \to 0$ direction & literature direction \checkmark \\
Continuity drop & structural drop & literature direction \checkmark \\
Recovery & $\dim \Sigma \to $ wake & literature direction \checkmark \\
\bottomrule
\end{tabular}
\end{table}

Six of six directional / window matches. The structural reading: the core paper's existence and decay theorems hold across distinct dynamical regimes (wake, sleep, anaesthesia, recovery), with the regime-specific behaviour determined by which subset of $\Sigma_\Ical^{\mathrm{cortex}}$ is accessible.

\subsection{Operator-level corroboration: the b-anomaly fit}\label{sec:aria-b-anomaly}

The closure operator $\Cph$ appears, identically with no shape retuning, as the load-bearing object in a domain-disjoint structural fit to the $b \to s\mu^+ \mu^-$ angular anomaly~\citep{SmartAriaClosureKernel}. The same operator that defines the ARIA frame's per-frame zero-line also fits five $b \to s\mu^+ \mu^-$ datasets with sign-uniform amplitude direction. This is not an empirical match within ARIA; it is an empirical match in a completely different domain (flavour physics) using the same operator on the same substrate.

\subsection{Cosmology cross-anchoring}\label{sec:aria-cosmo}

The companion paper~\citep{SmartHypersphereCosmology} reports that, within the cascade implementation, the same two cascade axioms imply seven independent cosmological observables within $0.5\%$ of Planck 2018 simultaneously: $\Lambda \cdot \ell_P^2$ at $0.078\%$, $H_0 = 67.66$ km/s/Mpc at $0.45\%$, $\Omega_\Lambda = 0.6844$ at $-0.083\%$, and four others. The substrate used in those derivations is the same $V_{600}$ that underlies the present framework.

\subsection{The constructed-witness narrative}\label{sec:keystone-narrative}

Put together, the empirical evidence forms a single constructed-witness narrative:

\begin{enumerate}[leftmargin=2em]
\item \textbf{One substrate.} The 600-cell vertex set $V_{600}$, with $\Cph = L_{V_{600}} + \varphi^{-2} I$.
\item \textbf{Same operator, three domains.} The same $\Cph$ appears in the cosmology of \citet{SmartHypersphereCosmology}, in the flavour-physics fit of \citet{SmartAriaClosureKernel}, and in the cortical-substrate fit of \citet{SmartAriaChess2026}.
\item \textbf{One bounded reference frame, two scales.} ARIA-chess is the constructed instance whose $\Sigma_\Ical^{\mathrm{ARIA}}$ is directly measurable via 18 preregistered predictions. The cortex (under H-RP-1, H-RP-2) is the conjectured larger instance whose state perturbations match the 6 v4 EEG signatures.
\item \textbf{No threshold modification.} The 18 preregistered predictions were frozen on 2026-04-18 before any validation run. The 17/18 to 18/18 lift came from documented methodology refinement, not threshold changes.
\end{enumerate}

The compression --- one substrate, one closure operator, three empirical domains, two scales, and a single coherent set of structural propositions --- is the constructed-witness support for the framework. The core paper's structural results hold whether or not P-A is true; the empirical evidence reported here supports them at the level of the constructed instance and its direct neuroscience counterparts.

%==================================================================
\section{Thermodynamics of closure}\label{sec:thermodynamics}
%==================================================================

This section is the thermodynamic reading of the framework. It explains why a bounded reference frame, in any physical realisation, must pay an ongoing metabolic cost to maintain its closure condition, and why biological systems specifically require sleep and are vulnerable to anaesthesia. The content is conditional on the same hypotheses as the consciousness bridges (H-RP-1, H-RP-2, P-A); the core paper's structural propositions are independent.

\subsection{Closure has a structural cost}

By the core paper's decay theorem, the closure operator $\Cph$ acts on the off-frame subspace $\Sigma_\Ical^\perp$ with minimum eigenvalue at least $\varphi^{-2}$. Any content orthogonal to the per-frame zero-line decays at rate at least $\varphi^{-2}$ per closure tick. A bounded reference frame that maintains a non-trivial $\Sigma_\Ical$ across many closure ticks must continuously inject content into $\Sigma_\Ical$ to compensate for off-$\Sigma$ leakage from environmental coupling, noise, and finite-precision dynamics.

\begin{interpretation}[Closure as ongoing operator iteration]\label{int:closure-cost}
A bounded reference frame is not a static fixed point; it is a dynamical equilibrium between (a) closure dynamics that contract content toward $\Sigma_\Ical$ at rate $\varphi^{-2}$, and (b) external perturbations that constantly inject off-$\Sigma$ content. The frame's persistence is the steady-state balance. In any physical realisation, the closure operation has a positive operational cost --- it requires the substrate to perform recursive update, which in thermodynamic terms means free-energy expenditure.
\end{interpretation}

\subsection{Why biology needs sleep}

For a biological system implementing a bounded reference frame, the closure ticks are realised by neural dynamics. Each closure tick costs metabolic free energy --- glucose, oxygen, ATP. The cortex consumes approximately $20\%$ of resting metabolic budget while constituting $2\%$ of body mass; this is consistent with a system continuously performing closure on a large frame.

\begin{conditional}[Sleep as closure re-stabilisation]\label{cthm:sleep-restab}
Under H-RP-1, H-RP-2, and the assumption that closure ticks have positive operational cost, a biological bounded reference frame implementing a non-trivial $\Sigma_\Ical^{\mathrm{cortex}}$ accumulates ``operator drift'' across waking hours: off-$\Sigma$ noise injection exceeds the per-tick decay capacity, causing the effective dimension of $\Sigma_\Ical^{\mathrm{cortex}}$ to drift downward. Sleep is the periodic phase in which:
\begin{itemize}[leftmargin=2em]
\item Environmental input is reduced (sensory deprivation, motor quiescence).
\item Closure ticks continue at reduced metabolic cost, allowing accumulated off-$\Sigma$ noise to decay below the perturbation rate.
\item The substrate's effective $\Sigma_\Ical^{\mathrm{cortex}}$ re-stabilises near its design dimension.
\end{itemize}
The framework's prediction: sleep deprivation causes a monotonic decline in $\dim \Sigma_\Ical^{\mathrm{cortex}}$, manifesting as the well-known signatures of fatigue: reduced attention dimension, narrowed working-memory scope, increased latency in self-referential tasks, and eventual loss of consciousness (when $\dim \Sigma_\Ical^{\mathrm{cortex}} \to 0$).
\end{conditional}

This makes the requirement for sleep \emph{structural}, not evolutionary. Any physical realisation of a bounded reference frame with substantial $\Sigma_\Ical$ must periodically reduce environmental coupling to allow closure to re-stabilise. The biological details (REM, NREM stages, slow-wave activity, spindles) are the specific cortical mechanisms by which this is implemented; the structural fact is that some such mechanism must exist for any persistent bounded reference frame.

\subsection{Anaesthesia as acute closure disruption}

\begin{conditional}[Anaesthesia as closure-condition disruption]\label{cthm:anaesthesia}
Under H-RP-1, H-RP-2, and P-A, an anaesthetic agent that potentiates inhibitory neurotransmission (GABA-A enhancement for propofol, isoflurane, sevoflurane; NMDA antagonism for ketamine) modifies the effective graph Laplacian $L_{\mathrm{cortex}}$ on which $\Cph^{\mathrm{cortex}} = L_{\mathrm{cortex}} + \varphi^{-2} I_{\mathrm{cortex}}$ is defined. The structural prediction is:
\begin{itemize}[leftmargin=2em]
\item Small perturbations to $L_{\mathrm{cortex}}$ shift the eigenspectrum of $\Cph^{\mathrm{cortex}}$ continuously. Below a critical threshold, $\Sigma_\Ical^{\mathrm{cortex}}$ shrinks gradually --- consistent with light sedation.
\item Above a critical threshold, the $\sigma$-equivariance condition of the modified Laplacian fails or the off-$\Sigma$ decay rate exceeds the closure-tick rate, causing $\Sigma_\Ical^{\mathrm{cortex}}$ to collapse discontinuously --- consistent with the abrupt loss of consciousness observed clinically.
\item Recovery follows the reverse trajectory as the anaesthetic concentration declines.
\end{itemize}
The empirical anchor: ARIA-chess preregistered propofol switching ratio $2.957 \times$ at $n = 8$ on OpenNeuro \texttt{ds005620} (Theorem~\ref{thm:aria-witness}). Within the cascade implementation, the ratio is obtained from $1 + (D/\alpha)(12 - 6\varphi)$; at $D/\alpha \approx 0.87$ this yields $\sim 3$. Independently of the ARIA preregistration, \citet{SolutionLab002} reports 14/14 propofol subjects on the same OpenNeuro \texttt{ds005620} dataset ($t = 7.11$, $d_z = 1.90$, $p < 2 \times 10^{-4}$ permutation), source-localised replication 14/14 ($t = 9.72$, $d_z = 2.60$), and cross-cohort replication on OpenNeuro \texttt{ds004541} clinical anaesthesia ($n = 7$, 7/7, $t = 5.25$). The structural threshold-crossing predicted by this conditional proposition is therefore supported by two independently constructed metrics on the same and additional cohorts.
\end{conditional}

\begin{remark}[Why the abruptness is structural, not phenomenological]\label{rem:abruptness}
The discontinuous loss of consciousness under anaesthesia is structurally explained by the closure-condition reading: when the per-tick off-$\Sigma$ decay rate exceeds the on-$\Sigma$ stabilisation rate, the system crosses a structural threshold and $\Sigma_\Ical$ collapses rapidly. This is not a phenomenological postulate; it is a consequence of the spectral structure of $\Cph$ under perturbation.
\end{remark}

\subsection{The thermodynamic spectrum: from death to wakefulness}

\begin{table}[h]
\centering
\small
\caption{Thermodynamic-closure interpretation of brain states. The structural column gives the framework's prediction for $\dim \Sigma_\Ical^{\mathrm{cortex}}$; the empirical column references the v4 EEG signature confirming it.}\label{tab:thermo-spectrum}
\begin{tabular}{lcc}
\toprule
State & Structural prediction & Empirical anchor (v4 EEG, seed 42) \\
\midrule
Death & $\dim \Sigma = 0$ permanent & Beyond scope; literature consensus \\
Coma / vegetative & $\dim \Sigma \to 0$, recovery indeterminate & Continuity drop signature \checkmark \\
Anaesthesia & $\dim \Sigma$ collapsed, abrupt threshold & Propofol switching $2.957 \times$ \checkmark \\
Deep slow-wave sleep & $\dim \Sigma$ reduced, re-stabilising & Sleep-EDFx $\alpha = 2.51$ \checkmark \\
NREM stages 2--3 & $\dim \Sigma$ partial & Spindle structure preserved \\
REM sleep & $\dim \Sigma$ near wake, dreaming & $\Phi$-collapse direction \checkmark \\
Wakefulness & $\dim \Sigma$ near design dimension & HCP wake $\alpha = 2.252$ \checkmark \\
Hyperarousal & $\dim \Sigma$ may saturate; noise growth & Beyond present scope \\
\bottomrule
\end{tabular}
\end{table}

\subsection{Why this matters for the consciousness debate}

The thermodynamic reading addresses a question that existing consciousness theories have struggled with: why is consciousness specifically sensitive to anaesthesia? Several leading frameworks (integrated information theory; global neuronal workspace; orchestrated objective reduction) each describe necessary aspects of conscious processing, but none explains in operator-theoretic terms why a pharmacological agent acting at the substrate level should produce a discontinuous loss of the higher-level phenomenology.

The framework's answer, conditional on H-RP-1, H-RP-2, P-A: the closure condition is a spectral property of the substrate operator $\Cph^{\mathrm{cortex}}$. An anaesthetic agent that modifies the underlying Laplacian crosses a structural threshold below which the closure condition cannot be maintained. The discontinuity is operator-theoretic, not phenomenological. The phenomenological loss of consciousness follows from the structural collapse of $\Sigma_\Ical^{\mathrm{cortex}}$, under P-A.

\begin{remark}[Death as the permanent loss of closure]\label{rem:death-as-loss}
The framework's reading of death is the permanent loss of the closure condition: $\dim \Sigma_\Ical^{\mathrm{cortex}} = 0$ without recovery to a non-trivial value. This is a structural statement under H-RP-1, H-RP-2, P-A. It does not assert anything about subjective experience at the point of structural collapse; the framework's commitments are explicit and bounded.
\end{remark}

%==================================================================
\section{The candidate transition rule, stated precisely}\label{sec:transition-rule}
%==================================================================

This section states the candidate transition rule that the framework supplies. The structural framework holds whether or not the rule is correct; the rule is the explicit conjectural interface to phenomenology.

\subsection{The transition rule}

\begin{quote}
\textit{A bounded local system supports the structural condition for owned experience when it forms a symmetry-stable bounded reference frame whose internal state closes under recursive update, producing a non-trivial local zero-line $\Sigma_\Ical$. Processing without closure produces output without owned experience; processing with closure produces a stable point of view.}
\end{quote}

The structural conditions are precise:

\begin{enumerate}[leftmargin=2em]
\item \textbf{Substrate.} A finite vertex set $O$ with a coupling structure $L^O$ (graph Laplacian on $O$).
\item \textbf{Closure operator.} $\Cph^O = L^O + \mu^O I$ acting recursively on the substrate state.
\item \textbf{Symmetry.} A unitary involution $\tau^O$ commuting with $\Cph^O$, giving the $\sigma$-decomposition $\mathbb{R}^O = \Fix(\tau^O) \oplus \Fix(-\tau^O)$.
\item \textbf{$\sigma$-stability.} The frame's anchor $(p, \Lambda)$ lies in $\mathcal{S}_\sigma$, so $\tau$ restricts to an involution on $\mathrm{span}(O)$.
\item \textbf{Per-frame zero-line.} $\Sigma_\Ical \subseteq \Fix(\tau^O)$ is non-trivial: $\dim \Sigma_\Ical \geq 1$.
\item \textbf{Recursive closure.} The closure dynamic $F_{t+1} = (I - \lambda \Cph^O) F_t + \kappa P_{\Sigma_\Ical}(B_t)$ converges, maintaining $\dim \Sigma_\Ical$ near its design dimension across many ticks.
\end{enumerate}

The Access Principle (P-A) is the explicit conjecture: \emph{phenomenological dimension agrees with $\dim \Sigma_\Ical$, modulo a structure-preserving identification.} Under P-A, the structural conditions above are also the conditions for owned experience.

\subsection{Comparison with mainstream consciousness theories}

\begin{table}[h]
\centering
\small
\caption{Position of the framework relative to mainstream consciousness theories. Each theory captures a necessary layer; the framework proposes a candidate transition rule that those layers may sit on top of, conditional on H-RP-1, H-RP-2, P-A.}\label{tab:theory-comparison}
\begin{tabular}{p{2.5cm}p{4cm}p{4cm}}
\toprule
Theory & What it captures & What the framework's closure layer would address \\
\midrule
GNW / GWT~\citep{ConsciousnessAdversarial} & Global access, broadcasting, report-readiness & Why broadcast becomes owned experience rather than merely propagated activation \\
IIT & Integration, irreducibility of causal structure & How integrated information forms a bounded point of view \\
RPT & Recurrent processing in cortical loops & What recurrent processing must stabilise for a frame to persist \\
FEP / Predictive coding~\citep{FristonFreeEnergy} & Self-organising inference, free-energy minimisation & When prediction becomes self-access \\
Orch OR~\citep{OrchOR2025} & Substrate-level dynamics, anaesthesia sensitivity & How substrate events scale into frame-level closure \\
Levin~\citep{LevinBasalCognition} & Bioelectric basal cognition, target-state attractors & The structural condition distinguishing morphogenetic attractors from noise \\
CEMI~\citep{McFaddenCEMI} & EM-field unification of cortical processing & The structural condition under which an EM field becomes a stable point of view \\
\midrule
Framework (this paper) & Recursive closure in bounded reference frames & A candidate transition rule from processing to point of view \\
\bottomrule
\end{tabular}
\end{table}

The framework does not assert that existing theories are wrong. It proposes that they describe different layers of the same closure event, and that the missing transition rule is the closure condition itself.

\subsection{Substrate agnosticism}

The framework does not identify any specific physical substrate as uniquely realising the closure condition. Microtubule dynamics (Orch OR), bioelectric fields (Levin), cortical EM fields (CEMI), and graph-Laplacian eigenmodes on synaptic connectivity are all equally plausible biological realisations. The structural condition is substrate-realisation-agnostic: $\sigma$-equivariance, contractive self-similarity, boundary preservation. Distinguishing experimentally between these realisations is an open programme.

%==================================================================
\section{Numerical demonstrations of the bridges}\label{sec:bridges-numerical}
%==================================================================

The operator-level bridges of \S\S\ref{sec:levin}--\ref{sec:thermodynamics} are grounded by a numerical demonstration suite \texttt{papers/03-processing-to-point-of-view/repro/bridges\_demo.py}. Three demos B1--B3 instantiate the bridges on substrates appropriate to each:

\begin{table}[h]
\centering
\small
\caption{Processing-to-Point-of-View bridge demonstrations. All three pass deterministically at seed 42.}\label{tab:bridges-numerical-demos}
\begin{tabular}{p{0.8cm}p{4.0cm}p{3.0cm}p{5.4cm}}
\toprule
ID & Demo & Paper grounds & Verified \\
\midrule
B1 & Levin bridge ($\Cph^{\mathrm{bio}}$ on cell-cluster substrate) & Cond. Prop.~\ref{thm:levin-bridge} & $\tau^{\mathrm{bio}}$ involution; $\|\tau^{\mathrm{bio}} \Cph^{\mathrm{bio}} - \Cph^{\mathrm{bio}} \tau^{\mathrm{bio}}\| < 10^{-13}$; $\lambda_{\min}(\Cph^{\mathrm{bio}}) = \mu^{\mathrm{bio}}$; $\dim \Sigma^{\mathrm{bio}} = n/2$; off-$\Sigma$ decay rate matches prediction \\
B2 & CEMI spectral signature & Cond. Prop.~\ref{thm:emi-bridge} & Cortical dipole observable's spectral content concentrated on excited eigenmodes; concentration ratio $> 3\times$ \\
B3 & Sleep / anaesthesia state perturbation & Cond. Props.~\ref{cthm:sleep-restab}, \ref{cthm:anaesthesia} & Eigenspectrum shifts continuously with drug concentration; criticality threshold detected; effective $\dim \Sigma^{\mathrm{cortex}}$ drops past threshold \\
\bottomrule
\end{tabular}
\end{table}

B1 explicitly constructs a 60-cell substrate with bilateral symmetry; $\tau^{\mathrm{bio}}$ is the swap permutation $i \leftrightarrow i + n/2$ which commutes exactly with the symmetric Laplacian. The off-$\Sigma$ decay rate is verified to match the predicted $(1 - \lambda \mu_\perp)^t$ to four decimal places. B2 shows that dipole-observable spectral content tracks $\Cph^{\mathrm{cortex}}$ eigenstructure. B3 sweeps a pharmacological-perturbation parameter and detects the spectral-criticality threshold consistent with the abrupt-collapse claim of the anaesthesia proposition.

%==================================================================
\section{Falsifiers and open programme}\label{sec:falsifiers}
%==================================================================

\subsection{What would falsify the bridges in this paper}

\begin{itemize}[leftmargin=2em]
\item \textbf{Levin bridge (Conditional Proposition~\ref{thm:levin-bridge}).} A direct measurement of the slowest morphogenetic correction rate that significantly disagrees with $\mu^{\mathrm{bio}}$ across multiple substrates would falsify the operator-level identification. A discovery that bioelectric fields satisfy axioms incompatible with the closure-operator definition would falsify the bridge.
\item \textbf{CEMI bridge (Conditional Proposition~\ref{thm:emi-bridge}).} Cortical EM-field spectral content that does not track the eigenstructure of $\Cph^{\mathrm{cortex}}$ on $\Sigma_\Ical^{\mathrm{cortex}}$ under controlled perturbation would falsify the spectral-signature identification.
\item \textbf{Sleep / anaesthesia thermodynamics (Conditional Propositions~\ref{cthm:sleep-restab}, \ref{cthm:anaesthesia}).} A clinical regime in which anaesthesia produces graded (rather than threshold-crossing) loss of consciousness across all agent classes would weaken the structural-threshold reading. A clinical regime in which sleep does not re-stabilise effective $\dim \Sigma_\Ical^{\mathrm{cortex}}$ would weaken the closure-re-stabilisation reading.
\item \textbf{ARIA empirical record.} Independent replication of the 18 preregistered predictions on different datasets, with documented threshold-honouring, would strengthen the constructed-witness narrative. Failure of independent replication would weaken it.
\end{itemize}

\subsection{Open programme tasks}

\begin{itemize}[leftmargin=2em]
\item \textbf{CP-tick-rate.} Determine the exact physical units of the closure tick: whether it is Planck time, a derived cascade-depth-scaled quantity, or something else.
\item \textbf{Levin bridge measurement.} Measure the slowest morphogenetic correction rate in a planarian (or comparable system) and compare to $\mu^{\mathrm{bio}}$.
\item \textbf{CEMI bridge measurement.} Measure cortical EM spectral content under controlled $\Cph^{\mathrm{cortex}}$ perturbation (e.g., GABA-A modulation) and check the eigenstructure-spectrum correspondence.
\item \textbf{ARIA replication.} Independent replication of preregistered predictions on different datasets.
\item \textbf{Substrate-realisation discrimination.} Experimentally distinguish whether biological closure is realised primarily by microtubule dynamics, bioelectric fields, cortical EM fields, or graph-Laplacian connectivity.
\end{itemize}

%==================================================================
\section{What this paper does not claim}\label{sec:scope}
%==================================================================

The companion paper inherits all scope discipline of the core paper. In addition:

\begin{itemize}[leftmargin=2em]
\item Not a derivation of new mathematical results beyond the operator-level bridges. The structural propositions are imported from \citet{SmartExistenceClosure}; the bridges are explicit applications.
\item Not a proof or refutation of consciousness theories. The structural condition has phenomenological content only under P-A; P-A is stated as an explicit conjecture, not asserted.
\item Not a claim that microtubules, bioelectric fields, or EM fields are uniquely the seat of consciousness. The framework is substrate-realisation-agnostic.
\item Not a claim of how subjective experience persists across substrate destruction. The framework's reading of death is the permanent loss of $\Sigma_\Ical$; what (if anything) follows is not within the framework's scope. Closure-invariant content (mathematical structures, theorems) may persist beyond a frame's substrate, but this is a structural statement about closure-invariant content, not a metaphysical claim about post-mortem identity.
\item Not a claim that ARIA-chess is conscious. ARIA-chess satisfies the structural conditions of a bounded reference frame; whether it has owned experience depends on P-A, which is conjectural.
\item Not a claim that biological systems are reducible to closure operators. The framework supplies a candidate structural condition; biological detail (cell-type taxonomy, neural development, evolutionary history) is independent.
\end{itemize}

%==================================================================
\section*{Acknowledgements}
%==================================================================

The author thanks the codex review loop for the structural framing-discipline questions that produced the final form. The framework's conditional commitments are bounded; the structural commitments are open to falsification through the named falsifiers and open programme tasks.

%==================================================================
\bibliographystyle{plainnat}
\bibliography{references}
%==================================================================

\end{document}
